\documentclass[10pt,a4paper]{amsart}
\usepackage[margin=19mm]{geometry}
\usepackage{amsmath,amssymb,mathtools}
\usepackage[scr=rsfs,cal=euler]{mathalfa}
\usepackage{xfrac}
\usepackage[unicode,hidelinks,bookmarks=false]{hyperref}
\newcommand{\ie}{i.e.\ }
\newcommand{\cf}{cf.\ }
\newcommand{\resp}{respectively\ }
\newcommand{\Subsection}{\subsection}
\providecommand{\nobreakdash}{\nobreak\hbox{-}}
\setlength{\emergencystretch}{2em}
\multlinegap0pt
\usepackage{amssymb}
\newtheorem{lem}{Lemma}
\newcommand{\e}{\varepsilon}
\newcommand{\til}{\widetilde}
\title{A tautological continuous field of Roe bimodules}
\author{Vladimir Manuilov}
\date{Source: arXiv:2603.23366v1 (CC BY 4.0)}
\begin{document}
\maketitle

\noindent\emph{Source and licence.} A tautological continuous field of Roe bimodules. Version arXiv:2603.23366v1 (CC BY 4.0), distributed by arXiv under the Creative Commons Attribution 4.0 International licence, http://creativecommons.org/licenses/by/4.0/, as recorded in the arXiv licence metadata for this exact version and shown on the abstract page https://arxiv.org/abs/2603.23366v1. Full licence notice: \texttt{licenses/arxiv-2603.23366-CC-BY-4.0.txt}.\par\smallskip

\noindent\textit{Editorial scope.} Counted unit: the article's lem base (source0.tex lines 546--556). The article's complete original proof of it is delivered with it (source0.tex lines 557--569, one \texttt{proof} environment of 13 lines). No mathematics is written, repaired or re-derived: the statement and the proof are the article's own lines, in the article's own notation, and no retained line is substituted or re-flowed.  No further source-local context is needed: every cross-reference of every retained line resolves inside the two delivered spans.  Nothing was omitted from any retained span, so the package carries no omission record.  No retained line contains a citation, so the wrapper carries no bibliography and no bibliography entry.  No retained line contains a figure environment, an image-inclusion command or drawing code, so no image is delivered and the package declares no asset dependency.  Editorial formatting only, in addition: an AMS \texttt{amsart} preamble that restates the article's own declarations and macros used by the retained lines, namely the environments \texttt{lem} on separate counters, together with the standard packages those lines need.  The retained environments print in this document as Lemma 1 in order of appearance, whereas the article prints the same items with the numbering of its own full document.  The complete native source of the article as distributed in the arXiv e-print for this exact version is bundled unchanged as \texttt{sources/197079/source0.tex}.
\begin{lem}
The collection of subsets 
\begin{equation}\label{base}
W_{A,B}=\{t\in T:\til\chi_{a_1}(t)=1,\ldots,\til\chi_{a_m}(t)=1;\til\chi'_{b_1}(t)=1,\ldots,\til\chi'_{b_n}(t)=1\},
\end{equation}
$$
A=\{a_1,\ldots,a_m\}\subset P,\quad B=\{b_1,\ldots,b_n\}\subset P.
$$ 
is a base for the Gelfand topology on $T=\gamma P$.  

\end{lem}

\begin{proof}
Recall that the Gelfand topology on $T=\gamma P$ is given by the subbase of sets 
\begin{equation}\label{subbase}
O_{s,f,\e}=\{t\in T:|\til f(t)-\til f(s)|<\e\}, \quad s\in T, f\in C, \e>0.
\end{equation}
Taking $f=\til\chi_{a_i}$ (resp., $f=\til\chi'_{b_j}$), $s\in U_{a_i}$ (resp., $s\in V_{b_j}$), $\e=1/2$, we get 
$$
O_{s,f,\e}=\{t\in T:|\til\chi_{a_i}(t)-1|<1/2\}=W_{(\{a_i\},\emptyset)}
$$ 
(resp. $O_{s,f,\e}=W_{(\emptyset,\{b_j\})}$), hence the topology determined on $T$ by the subbase (\ref{subbase}) is finer than the topology determined by the base (\ref{base}). As $P$ is dense in the Gelfand topology, this means that $P$ is dense in the topology defined by (\ref{base}) as well.  

To see the opposite inclusion, note that density of $P$ implies that the subbase (\ref{subbase}) can be made smaller by using only sets $O_{s,f,\e}$ with $s\in P$, and that both topologies coincide on $P$. For $f\in C_b(P)$ and $O_{s,f,\e}$ there exists, by continuity of $f$, an open set $V\subset P$ such that $V\subset O_{s,f,\e}\cap P$. As $V$ is open, there exists a basic set $U\subset P$ such that $U\subset V$. Then there exists $W_{A,B}$ such that $W_{A,B}\cap P=U$. Thus, $W_{A,B}\subset O_{s,f,\e}$, so the topology on $T$ with the base (\ref{base}) is finer than the Gelfand topology.  
\end{proof}

\end{document}
